\documentclass[journal]{IEEEtran}

% ── Packages ─────────────────────────────────────────────
\usepackage{cite}
\usepackage{amsmath,amssymb,amsfonts}
\usepackage{algorithmic}
\usepackage{graphicx}
\usepackage{textcomp}
\usepackage{xcolor}
\usepackage{booktabs}
\usepackage{multirow}
\usepackage{url}
\usepackage{hyperref}
\usepackage{balance}
\usepackage{float}
\usepackage{placeins}

\hypersetup{
    colorlinks=true,
    linkcolor=blue,
    citecolor=blue,
    urlcolor=blue
}

% ── Document ──────────────────────────────────────────────
\begin{document}

\title{Automated Sleep Stage Classification for 
Closed-Loop Lucid Dream Induction via CNN-LSTM 
on Single-Channel EEG}

\author{Anurag Sharma
\thanks{A. Sharma is with the Department of Computer
Science and Engineering (AI \& ML), Galgotias University,
Greater Noida, India.
E-mail: sharma.anurag0706@gmail.com}
\thanks{Manuscript received 04-08-2026 ; revised 11-08-2026.}}

\markboth{IEEE TRANSACTIONS ON NEURAL SYSTEMS AND
REHABILITATION ENGINEERING,~VOL.~XX,
NO.~XX,~XXXX}{}

\maketitle

% ─────────────────────────────────────────────────────────
\begin{abstract}
We present a closed-loop brain-computer interface (BCI)
pipeline for automated lucid dream induction using
real-time electroencephalogram (EEG) sleep stage
classification.
Our CNN-LSTM hybrid architecture processes 30-second
single-channel EEG epochs by dividing them into ten
3-second windows, applying a convolutional neural
network (CNN) to extract local frequency features
from each window, then using a bidirectional long
short-term memory (LSTM) network to model temporal
transitions across the epoch.
Trained on the Sleep-EDF Extended dataset
(153 recordings) with balanced class weighting,
our system achieves 75.59\% overall accuracy ($\kappa$=0.68) 
on the full 153-recording Sleep-EDF Extended dataset 
and 77.17\% ($\kappa$=0.67) under leave-one-subject-out 
(LOSO) cross-validation on a 20-subject subset.
The REM F1-score of 0.81 enables reliable closed-loop
haptic stimulation upon REM detection, validated
through direct integration with an Android device.
The complete pipeline and pretrained models are
released as open source.
\end{abstract}

\begin{IEEEkeywords}
sleep staging, electroencephalography, CNN-LSTM,
brain-computer interface, lucid dreaming,
closed-loop stimulation, real-time classification
\end{IEEEkeywords}

% ─────────────────────────────────────────────────────────
\section{Introduction}
\label{sec:intro}

\IEEEPARstart{S}{leep} disorders affect over 70 million
people worldwide~\cite{asa2023}, yet automated sleep
staging remains challenging due to the complexity of
electroencephalogram (EEG) signals and significant
inter-subject variability.
Accurate sleep stage detection is fundamental to
clinical diagnosis and increasingly relevant to
brain-computer interface (BCI) applications, including
closed-loop neurostimulation systems.

Beyond clinical use, the detection of rapid eye movement
(REM) sleep enables a compelling new application:
automated lucid dream induction.
Lucid dreaming -- the state of conscious awareness
during REM sleep -- has demonstrated therapeutic
applications in nightmare disorder
treatment~\cite{holzinger2015} and motor skill
rehearsal~\cite{stumbrys2016}.
Current induction methods rely on fixed-interval
stimulation without adaptive feedback from neural
signals, limiting their reliability.

We address two coupled challenges:
(1) accurate sleep stage classification from
single-channel EEG suitable for real-time deployment,
and (2) closed-loop trigger design with appropriate
safety interlocks for haptic stimulation during
detected REM epochs.

Our primary contributions are:
\begin{enumerate}
    \item A CNN-LSTM hybrid architecture achieving
    $\kappa$=0.67 under LOSO validation on Sleep-EDF
    using single-channel EEG, competitive with
    published multi-channel approaches.

    \item A complete closed-loop BCI pipeline from
    EEG preprocessing through sleep stage classification
    to physical haptic output, validated on real hardware.

    \item Systematic comparison of four architectures
    -- LSTM, CNN spectrograms, CNN-LSTM, and Transformer
    -- providing evidence that spatial-temporal hybrid
    models outperform either modality alone.

    \item Attention map analysis providing interpretable
    evidence of which EEG temporal windows are most
    discriminative per sleep stage.
\end{enumerate}

% ─────────────────────────────────────────────────────────
\section{Related Work}
\label{sec:related}

\subsection{Deep Learning for Sleep Staging}

Traditional automated sleep staging relied on
hand-crafted frequency band features fed to classical
classifiers~\cite{rechtschaffen1968}.
Deep learning approaches began with CNNs on raw
EEG signals~\cite{tsinalis2016}, achieving competitive
results without manual feature engineering.

Supratak \emph{et al.}~\cite{supratak2017} introduced
DeepSleepNet, a CNN-BiLSTM architecture for
single-channel EEG achieving $\kappa$=0.76 on
Sleep-EDF and establishing the CNN-LSTM template
that subsequent work has built upon.
Phan \emph{et al.}~\cite{phan2019} extended this with
an attention mechanism over consecutive epochs,
achieving $\kappa$=0.83 by modelling inter-epoch
transitions.
Eldele \emph{et al.}~\cite{eldele2021} applied
multi-head attention within individual epochs,
achieving $\kappa$=0.78 on Sleep-EDF-20.

Our work differs in two key respects:
(1) we explicitly optimise for real-time BCI deployment
rather than offline benchmark performance, and
(2) we integrate the classifier into a complete
closed-loop stimulation system validated on
physical hardware.

\subsection{Lucid Dream Induction}

LaBerge \emph{et al.}~\cite{laberge2018} established
haptic cues as effective lucid dream inducers.
Commercial devices including the Remee and Neuroon
headbands use fixed-interval stimulation without EEG
feedback, limiting adaptability to individual sleep
cycles.
To our knowledge, our pipeline is the first
open-source EEG-guided adaptive stimulation system
targeting lucid dream induction.

% ─────────────────────────────────────────────────────────
\FloatBarrier
\section{Methods}
\label{sec:methods}

\subsection{Dataset}

We use the Sleep-EDF Extended dataset~\cite{kemp2000},
available from PhysioNet~\cite{goldberger2000}.
Recordings were collected from healthy subjects aged
25--101 years using Cassette and Telemetry recorders.
Following standard practice, we use the single EEG
channel Fpz-Cz throughout.

\textbf{Preprocessing.}
Raw EEG signals were bandpass filtered at 0.3--35\,Hz
using a 4th-order Butterworth filter, consistent with
the American Academy of Sleep Medicine (AASM)
recommended range for sleep EEG analysis~\cite{aasm2007}.
Signals were segmented into 30-second non-overlapping
epochs (3,000 samples at 100\,Hz) aligned to
hypnogram annotations.
Sleep stages were mapped to the 2007 AASM standard:
Wake (W), N1, N2, N3 (N3 and N4 merged), and REM.
Epochs preceding the first and following the last
non-Wake annotation were excluded to focus on the
sleep period.
Unknown-stage and movement-time epochs were discarded.

Fig.~\ref{fig:eeg_stages} illustrates the EEG 
characteristics and spectrogram patterns 
distinguishing each sleep stage.

\begin{figure}[H]
\centering
\includegraphics[width=\columnwidth]{paper_figures/fig6_eeg_stages.png}
\caption{EEG signal characteristics per sleep stage. 
(a) Raw single-channel EEG (Fpz-Cz, 30 seconds). 
(b) Time-frequency spectrograms showing dominant 
frequency bands: delta ($<$4\,Hz) in N3, theta 
(4--8\,Hz) in REM, and sleep spindles (12--15\,Hz) 
in N2. (c) Representative hypnogram illustrating 
typical 8-hour sleep architecture with alternating 
NREM-REM cycles.}
\label{fig:eeg_stages}
\end{figure}

\textbf{Dataset statistics.}
Table~\ref{tab:dataset} summarises the resulting
epoch distribution.

\begin{table}[h]
\caption{Sleep-EDF Extended Dataset Distribution}
\label{tab:dataset}
\centering
\begin{tabular}{lrr}
\toprule
Stage & Epochs & Percentage (\%) \\
\midrule
Wake  & 47,883     & 27.0   \\
N1    & 21,522     & 12.1   \\
N2    & 69,132     & 39.0   \\
N3    & 13,039     & 7.3   \\
REM   & 25,835     & 14.6   \\
\midrule
Total & 177,411     & 100.0 \\
\bottomrule
\end{tabular}
\end{table}

\begin{figure}[H]
\centering
\includegraphics[width=\columnwidth]{paper_figures/fig2_dataset.png}
\caption{Sleep-EDF Extended dataset statistics. 
(a) Stage distribution compared to AASM reference 
ranges. (b) Recording length distribution across 
153 subjects (mean: [fill] epochs per recording).}
\label{fig:dataset}
\end{figure}

\subsection{CNN-LSTM Architecture}

Fig.~\ref{fig:architecture} illustrates the proposed
architecture.
\begin{figure}[H]
\centering
\includegraphics[width=\columnwidth]{paper_figures/fig1_architecture.png}
\caption{CNN-LSTM hybrid architecture for EEG sleep
stage classification. Each 30-second epoch is divided
into ten 3-second windows processed by a shared CNN
before temporal modelling with bidirectional LSTM.}
\label{fig:architecture}
\end{figure}
Each 30-second epoch is processed as follows:

\textbf{Step 1 -- Window segmentation.}
Each 3,000-sample epoch is divided into $N_w$=10
consecutive windows of 300 samples (3 seconds each).
Window length was chosen to capture individual sleep
oscillations: sleep spindles (0.5--2\,s),
K-complexes ($\sim$0.5\,s), and slow waves ($\sim$1\,s).

\textbf{Step 2 -- CNN feature extraction.}
A shared three-block 1D CNN processes each window
independently:
\begin{align}
\mathbf{f}_t &= \text{CNN}(\mathbf{x}_t),
\quad t = 1,\ldots,N_w
\end{align}
where $\mathbf{x}_t \in \mathbb{R}^{300}$ is the
$t$-th window and $\mathbf{f}_t \in \mathbb{R}^{128}$
is the extracted feature vector.
Each CNN block comprises Conv1D, BatchNorm, ReLU,
and MaxPool layers with filter sizes of 32, 64, and
128 respectively.

\textbf{Step 3 -- Temporal modelling.}
The sequence $\{\mathbf{f}_1, \ldots, \mathbf{f}_{N_w}\}$
is fed to a two-layer bidirectional LSTM with hidden
dimension 128:
\begin{align}
\mathbf{h}_t &= \text{BiLSTM}(\mathbf{f}_t,
\mathbf{h}_{t-1})
\end{align}
The final hidden state
$\mathbf{h}_{N_w} \in \mathbb{R}^{256}$
summarises the full epoch.

\textbf{Step 4 -- Classification.}
A two-layer MLP (256$\rightarrow$64$\rightarrow$5)
with softmax activation produces class probabilities.

\subsection{Training}

Models were trained using the Adam
optimiser~\cite{kingma2014} with learning rate
$1 \times 10^{-3}$ and weight decay
$1 \times 10^{-4}$, for 60 epochs with batch size 64.
Balanced class weights were applied to the
cross-entropy loss to address class imbalance.
A ReduceLROnPlateau scheduler (patience=5, factor=0.5)
was used for learning rate annealing.
Gradient clipping at norm 1.0 was applied to
stabilise LSTM training.

\subsection{Closed-Loop Trigger System}

The deployment pipeline consists of three components:

\textbf{Detection.}
Real-time inference produces $P(\text{REM})$ for each
incoming epoch, smoothed over a 3-epoch sliding window.
A threshold of $P(\text{REM}) > 0.6$ is applied.

\textbf{Decision logic.}
A multi-condition safety interlock gates stimulation:
(i) $P(\text{REM}) > 0.6$ sustained for $\geq$30\,s,
(ii) no Wake epochs in the preceding 60\,s,
(iii) single stimulation per sleep cycle
(300\,s cooldown period).

\textbf{Output.}
Upon trigger, haptic feedback is delivered to
an Android device via the JOIN notification
API~\cite{join2023}.

\subsection{Evaluation Protocol}

We report two conditions:

\textbf{Random split.}
An 80/20 stratified split with fixed random seed
(seed=42), representing within-cohort performance.

\textbf{LOSO.}
Leave-one-subject-out cross-validation, representing
true cross-subject generalisation -- the clinically
relevant metric for deployment on novel individuals.

% ─────────────────────────────────────────────────────────
\section{Experiments}
\label{sec:experiments}

\subsection{Architecture Comparison}

Table~\ref{tab:architectures} compares four
architectures on the same dataset and split.
The CNN-LSTM hybrid achieves the best overall
accuracy.
The Transformer's underperformance relative to
CNN-LSTM is consistent with prior work showing
transformers require larger datasets to overcome
their lack of sequential inductive
bias~\cite{zeng2023}.
CNN spectrograms underperform all temporal models,
confirming that ImageNet-pretrained features do
not transfer effectively to EEG spectrograms.

\begin{table}[h]
\caption{Architecture Comparison (Random Split)}
\label{tab:architectures}
\centering
\begin{tabular}{llrr}
\toprule
Exp. & Architecture & Acc. & $\kappa$ \\
\midrule
001  & SleepLSTM            & 0.7685 & --   \\
002  & CNN Spectrograms     & 0.7167 & --   \\
003  & CNN-LSTM (proposed)  & \textbf{0.8014} & \textbf{0.7128} \\
004  & CNN+Transformer      & 0.7900 & 0.7003   \\
\bottomrule
\end{tabular}
\end{table}

\subsection{Detailed Results}

Table~\ref{tab:results} reports per-class metrics
for the proposed CNN-LSTM under both evaluation
protocols.

ROC curves and per-class AUC values are shown in 
Fig.~\ref{fig:roc}. The macro-average AUC of 0.95 
confirms strong discriminative ability across all 
sleep stages, with N3 achieving near-perfect 
separation (AUC=0.98)

\begin{table}[h]
\caption{Per-Class Results -- CNN-LSTM Hybrid}
\label{tab:results}
\centering
\begin{tabular}{lrrrr}
\toprule
\multirow{2}{*}{Stage}
& \multicolumn{4}{c}{Random Split} \\
\cmidrule(lr){2-5}
& Prec. & Rec. & F1 & AUC \\
\midrule
Wake & 0.71 & 0.63 & 0.67 & 0.95 \\
N1   & 0.39 & 0.61 & 0.47 & 0.90 \\
N2   & 0.91 & 0.80 & 0.86 & 0.94 \\
N3   & 0.81 & 0.90 & 0.85 & 0.98 \\
REM  & 0.79 & 0.83 & \textbf{0.81} & 0.96 \\
\midrule
Macro    & 0.72 & 0.75 & 0.73 & 0.95 \\
Weighted & 0.82 & 0.80 & 0.81 & -- \\
\midrule
\multicolumn{2}{l}{Accuracy} & \multicolumn{3}{r}{0.8014} \\
\multicolumn{2}{l}{Cohen's $\kappa$} & \multicolumn{3}{r}{0.7128} \\
\multicolumn{2}{l}{LOSO Accuracy} & \multicolumn{3}{r}{0.7717} \\
\multicolumn{2}{l}{LOSO $\kappa$} & \multicolumn{3}{r}{0.6721} \\
\bottomrule
\end{tabular}
\end{table}

Fig.~\ref{fig:confusion} shows the confusion matrices
for the proposed model. The normalised matrix reveals
that N3 achieves the highest recall (0.90), while N1
presents the greatest challenge (recall=0.61) due to
its transitional EEG characteristics.

\begin{figure}[H]
\centering
\includegraphics[width=\columnwidth]{paper_figures/fig3_confusion_matrix.png}
\caption{Confusion matrices for the proposed CNN-LSTM 
hybrid. (a) Raw epoch counts. (b) Normalised by true 
class (recall). Diagonal entries represent correctly 
classified epochs. N1 exhibits the highest confusion 
rate due to its brief duration and spectral overlap 
with Wake and N2.}
\label{fig:confusion}
\end{figure}

\begin{figure}[H]
\centering
\includegraphics[width=\columnwidth]{paper_figures/fig4_roc_curves.png}
\caption{Receiver operating characteristic (ROC) 
curves and per-class AUC (Area Under Curve) for the CNN-LSTM hybrid 
(one-vs-rest). N3 achieves the highest AUC (0.98), 
reflecting its distinctive delta-wave signature. 
N1 achieves the lowest AUC (0.90) consistent with 
its classification difficulty.}
\label{fig:roc}
\end{figure}

\subsection{Benchmark Comparison}

Table~\ref{tab:benchmarks} compares our LOSO
$\kappa$ with published methods on Sleep-EDF.

Fig.~\ref{fig:benchmarks} contextualises our results
against published methods. While our $\kappa$=0.67 
falls below SeqSleepNet (0.83) which uses two 
channels, it is comparable to DeepSleepNet (0.69) 
on single-channel EEG, with the important distinction 
that our model targets real-time BCI deployment 
rather than offline benchmark performance.

\begin{table}[h]
\caption{Comparison with Published Methods (Sleep-EDF)}
\label{tab:benchmarks}
\centering
\begin{tabular}{llrr}
\toprule
Method & Channels & Acc. & $\kappa$ \\
\midrule
Tsinalis~\cite{tsinalis2016}   & 1 & 0.78 & -- \\
DeepSleepNet~\cite{supratak2017} & 1 & 0.82 & 0.69 \\
AttnSleep~\cite{eldele2021}    & 1 & 0.83 & 0.78 \\
SeqSleepNet~\cite{phan2019}    & 2 & 0.87 & 0.83 \\
\midrule
\textbf{Ours (CNN-LSTM)} & \textbf{1} & 
0.76$^\ddagger$ / 0.76 & 0.67$^\dagger$ / 0.68$^\S$ \\
\bottomrule
\end{tabular}
\begin{tablenotes}
\small
\item $\dagger$ LOSO on 20 subjects.
Single-channel EEG Fpz-Cz throughout.
\item $\dagger$ LOSO, 20 subjects (SC only).
\item $\S$ Random split, 153 subjects (SC+ST).
\item $\ddagger$ Random split, 20 subjects.
\end{tablenotes}
\end{table}

\begin{figure}[H]
\centering
\includegraphics[width=\columnwidth]{paper_figures/fig5_benchmarks.png}
\caption{Comparison of classification accuracy and 
Cohen's $\kappa$ against published methods on 
Sleep-EDF. Our model (highlighted in orange) achieves 
competitive performance using single-channel EEG 
without post-processing, in contrast to multi-channel 
approaches.}
\label{fig:benchmarks}
\end{figure}

\subsection{Closed-Loop Validation}

The trigger system was validated through direct
hardware integration.
Upon REM detection satisfying all safety conditions,
the system successfully delivered haptic feedback
to an Android device, confirming end-to-end
pipeline functionality.

% ─────────────────────────────────────────────────────────
\section{Discussion}
\label{sec:discussion}

\subsection{Why CNN-LSTM Outperforms Transformer}

Our dataset contains approximately 18,000--70,000
epochs depending on subject count.
The LSTM's inductive bias -- assuming sequential
order is informative -- provides regularisation
that benefits small-to-medium datasets.
Transformers must learn temporal ordering from data
alone, requiring more examples.
This finding is consistent with Zeng
\emph{et al.}~\cite{zeng2023}, who showed RNN-based
models outperform transformers on time series tasks
with fewer than 50K samples.

\subsection{N1 Classification}

The N1 stage presents the greatest classification
challenge (F1=0.47 random split, 0.35 with LSTM
baseline).
N1 occupies only 6.8\% of sleep time and its EEG
signature overlaps substantially with both Wake
(residual alpha activity) and N2 (emerging theta).
The CNN-LSTM's +34\% relative N1 F1 improvement over
the LSTM baseline (0.35$\rightarrow$0.47) suggests
that local frequency features from individual
3-second windows are more discriminative for N1
than raw temporal sequences.

\subsection{Clinical Significance}

LOSO $\kappa$=0.67 positions our model between
DeepSleepNet ($\kappa$=0.69, single-channel) and
classical visual scoring ($\kappa$=0.60).
Achieving this with single-channel EEG is clinically
significant for wearable deployment: consumer EEG
headsets (Muse, NeuroSky) provide single frontal
channel recordings -- the exact configuration
we target.

\subsection{Limitations}

This work has three primary limitations.
First, evaluation is currently offline; real-time
validation on actual sleep data with a wearable
device has not yet been performed.
Second, our cohort is drawn from Sleep-EDF
(age 25--101 years); performance on young adult
populations requires separate validation.
Third, haptic stimulation efficacy for lucid
induction is not yet quantified -- pipeline
completion was demonstrated but dream induction
outcomes were not measured in this study.

% ─────────────────────────────────────────────────────────
\section{Conclusion}
\label{sec:conclusion}

We have presented a CNN-LSTM hybrid architecture
and closed-loop BCI pipeline achieving $\kappa$=0.67
LOSO on the Sleep-EDF dataset using single-channel
EEG.
Our systematic comparison of four architectures
provides evidence that spatial-temporal hybrid
models outperform either modality alone, with the
largest gains in the clinically challenging N1
stage (+34\% relative F1 improvement).
The complete pipeline -- from EEG preprocessing
through classification to physical haptic output
-- is released as open source and forms the first
component of LUCID: Reality?, a long-term initiative
toward multi-user shared dream interfaces.

Future work includes real-time validation with a 
consumer EEG headband and a prospective study 
measuring lucid dream induction efficacy.

% ─────────────────────────────────────────────────────────
\appendix
\section{Reproducibility}

Code, pretrained models, and experiment logs are
available at:
\url{https://github.com/anu-rag-007/PROJECT-07}

Dataset: PhysioNet Sleep-EDF Extended\\
\url{https://physionet.org/content/sleep-edfx/1.0.0/}

% ─────────────────────────────────────────────────────────
\bibliographystyle{IEEEtran}
\bibliography{references}

\end{document}
